\section{Methodology}
\label{sec:methodology}

Our methodology is motivated directly by the gradient dead zone problem. To prove that ratio reward addresses binary reward's zero-variance property in GRPO groups, we use a two-phase design: (1) a mechanistic existence proof (h-e1) that isolates the group-level gradient signal difference, and (2) a policy-level training experiment (h-m1) that tests whether the mechanistic difference produces downstream performance improvements.

\subsection{Reward Functions}

We compare two reward formulations for GRPO post-training of code LLMs.

\textbf{Binary reward.} Given a completion and $n$ test cases, binary reward assigns:
\begin{equation}
r_{\text{binary}}(c) = \mathbf{1}[\forall i \in [n]: \text{pass}(c, t_i)]
\end{equation}
where $\text{pass}(c, t_i) = 1$ if the completion $c$ produces the correct output on test case $t_i$ within a 5-second timeout, and 0 otherwise. This is the reward used by all prior RLEF work for code LLMs.

\textbf{Ratio reward.} Ratio reward assigns credit proportional to the fraction of test cases passed:
\begin{equation}
r_{\text{ratio}}(c) = \frac{1}{n} \sum_{i=1}^{n} \text{pass}(c, t_i) \in [0, 1]
\end{equation}
Ratio reward coincides with binary reward when all tests pass (both = 1) or none pass (both = 0). It differs precisely in the partial-pass regime ($0 < k < n$ tests passing).

\textbf{Implementation.} Both functions are implemented as subprocess sandboxes with a 5-second timeout per test case. The implementation is validated by 12/12 unit tests covering pass/fail/timeout/error scenarios.

\subsection{GRPO Advantage Computation and the Dead Zone}

GRPO computes advantages within a group of $G$ completions for a given prompt by:
\begin{equation}
A_i = r_i - \bar{r}, \quad \bar{r} = \frac{1}{G} \sum_{j=1}^{G} r_j
\end{equation}
The group gradient contribution is proportional to the variance of advantages. When all completions fail all test cases (the all-zero binary reward regime), $r_i = 0$ for all $i$, $\bar{r} = 0$, and $A_i = 0$ for all $i$ --- the group contributes zero gradient to any model parameter.

\textbf{Formal condition.} Binary reward produces zero advantage variance in group $G$ if and only if all $k_i \in \{0, n\}$ (all-fail or all-pass for every completion). Under binary reward, a group where no completion passes all $n$ tests but some completions pass $1 \leq k < n$ tests (partial credit) produces zero variance --- this is the dead zone. Ratio reward produces non-zero variance whenever $\exists i: 1 \leq k_i < n$ and completions differ in pass count, which is the complement of the dead zone condition.

\subsection{Phase 1: Mechanistic Existence Proof (h-e1)}

\textbf{Objective.} Prove that ratio reward provides non-zero advantage variance in GRPO groups where binary reward provides zero variance, and confirm that the training infrastructure is functionally correct.

We use three components: (1) direct synthetic group computation (the illustrative Table~\ref{tab:advantages} example), (2) a 1,000-group simulation under realistic early-training conditions ($G=8$, $n=5$ test cases, $p_{\text{pass}}=0.1$ per test), and (3) a 3-step smoke test on actual GRPO training.

\textbf{Gate criterion.} h-e1 is satisfied if: code executes without error, the mechanistic proof demonstrates advantage variance difference, and gradient norms are measurable.

\subsection{Phase 2: Policy-Level Mechanism Test (h-m1)}

\textbf{Objective.} Test whether the gradient signal difference translates to measurable policy-level performance differences.

\textbf{Setup.} Two conditions --- binary and ratio reward --- trained with GRPOTrainer (trl 1.0.0) for up to 1,000 steps on APPS training split. Model: DeepSeek-Coder-6.7B-instruct. Separate H100 NVL GPUs per condition. Hyperparameters: \texttt{group\_size=8}, \texttt{learning\_rate=1e-6}, \texttt{kl\_beta=0.04}, \texttt{max\_new\_tokens=512}, \texttt{seed=42}.

\textbf{Key diagnostic: FractionPartialCallback.} Monitors fraction of completions per batch with $0 < k < n$ test cases passing at every step. If \texttt{fraction\_partial = 0} consistently, ratio reward is degenerate (equivalent to binary) and the experiment is terminated with Phase 0 redesign routing.

\textbf{Gate criterion.} Ratio HumanEval pass@1 $-$ binary $\geq 0.03$ with 95\% bootstrap CI excluding 0.

\subsection{Training Configuration}

\begin{table}[h]
\centering
\caption{Training hyperparameters.}
\label{tab:config}
\begin{tabular}{ll}
\toprule
\textbf{Parameter} & \textbf{Value} \\
\midrule
Model & deepseek-ai/deepseek-coder-6.7b-instruct \\
Dataset & codeparrot/apps \\
Group size & 8 \\
Learning rate & $1 \times 10^{-6}$ \\
KL beta & 0.04 \\
Max new tokens & 512 (h-m1); 256 (h-e1 smoke test) \\
GPU & H100 NVL (95 GB each) \\
Framework & trl 1.0.0, torch 2.5+cu124, Python 3.10 \\
\bottomrule
\end{tabular}
\end{table}
